\subsection{单环的情形}

设 A 与 B 是单环，f 是从 A 到 B 的一个同态。设 S 是一个单 A-模，T 是一个单 B-模。我们令

\[
i (f) = [ f ^ {*} (\mathrm{S}): \mathrm{T} ] = \mathrm{long} _ {\mathrm{B}} (f ^ {*} (\mathrm{S}));\tag{26}
\]

我们把基数 \(i(f)\) 称为 f 的指数。当 A 是 B 的子环且 f 是 A 到 B 中的典范单射时，我们把 \(i(\mathrm{B},\mathrm{A})\) 记作 \(i(f)\) 的替代，并把这个基数称为 A 在 B 中的指数。我们类似地定义 f 的高度 \(h(f)\)：

\[
h (f) = [ f _ {*} (\mathrm{T}): \mathrm{S} ] = \mathrm{long} _ {\mathrm{A}} (f _ {*} (\mathrm{T})).\tag{27}
\]

当 A 是 B 的子环且 f 是 A 到 B 中的典范单射时，我们用 \(h(\mathrm{B}, \mathrm{A})\) 表示 \(h(f)\)，并把它称为 A 在 B 中的高度。

A-模 S 是单生成的，故 B-模 \(f^{*}(\mathrm{S}) = \mathrm{B} \otimes_{\mathrm{A}} \mathrm{S}\) 亦然。由此得出 \(i(f)\) 是有限的，从而是一个整数。设 M 是一个 A-模。用 a 表示它的长度；则 M 同构于 \(\mathrm{S}^{(\mathfrak{a})}\)。于是 B-模 \(f_{*}(\mathrm{M})\) 同构于 \(f_{*}(\mathrm{S})^{(\mathfrak{a})}\)。由 \(i(f)\) 的定义，我们因此有

\[
\mathrm{long} _ {\mathrm{B}} (f ^ {*} (\mathrm{M})) = i (f) \mathrm{long} _ {\mathrm{A}} (\mathrm{M}).\tag{28}
\]

Z-模 \(K_{0}(A)\) 与 \(K_{0}(B)\) 都是维数为 1 的自由模，分别以 {[}S{]} 与 {[}T{]} 为基，并且我们有

\[
f ^ {*} ([ \mathrm{S} ]) = i (f) [ \mathrm{T} ].\tag{29}
\]

特别地，取 \(M = A_{s}\)；则 \(f^{*}(A_{s}) = B \otimes_{A} A_{s}\) 同构于 \(B_{s}\)（II, §3, No.~4, 第 249 页），故

\[
i (f) = \mathrm{long(B)} / \mathrm{long(A)}.\tag{30}
\]

由韦德伯恩定理（VIII, 第 120 页，定理 1），存在整数 \(m \geqslant 1\) 与 \(n \geqslant 1\) 以及体 D 与 E，使得 A 同构于 \(\mathbf{M}_m(\mathrm{D})\)，B 同构于 \(\mathbf{M}_n(\mathrm{E})\)。由公式 (30)，我们有 \(i(f) = n / m\)；特别地，\(m\) 整除 \(n\)。

设 N 是一个 B-模；用 a 表示它的长度。则 N 同构于 \(\mathrm{T}^{(\mathfrak{a})}\)，故 A-模 \(f_{*}(\mathrm{N})\) 同构于 \(f_{*}(\mathrm{T})^{(\mathfrak{a})}\)；由 \(h(f)\) 的定义，我们有

\[
\mathrm{long} _ {\mathrm{A}} (f _ {*} (\mathrm{N})) = h (f) \mathrm{long} _ {\mathrm{B}} (\mathrm{N}).\tag{31}
\]

我们已经看到（VIII, 第 124 页，命题 5），f 使 B 成为一个自由 A-模，且该模的所有基都具有同一基数，记作 \([B:A]_{s}\)，称为 B 在 A 上的（左）次数。A-模 \(f_{*}(B_{s})\) 同构于 \(A_{s}^{[B:A]_{s}}\)，故其长度为 \([B:A]_{s}\) long(A)。把公式 (30) 与公式 (31) 应用于特例 \(N=B_{s}\)，我们因此有

\[
[ \mathrm{B}: \mathrm{A} ] _ {s} = i (f) h (f).\tag{32}
\]

现在假设 B 是一个有限生成 A-模，即 \([B:A]_{s}\) 是有限的。则由公式 (32)，\(h(f)\) 是有限的。我们（在 VIII, 第 202 页）定义了一个群同态 \(f_{*}\)（从 \(K_{0}(B)\) 到 \(K_{0}(A)\)）；我们有

\[
f _ {*} ([ \mathrm{T} ]) = h (f) [ \mathrm{S} ].\tag{33}
\]

假设 A 与 B 是交换域 K 上为有限维的代数，且 f 是 K-线性的。同前，存在整数 \(m \geqslant 1\) 与 \(n \geqslant 1\)、为体的 K-代数 D 与 E，以及从 A 到 \(\mathbf{M}_{m}(\mathrm{D})\)、从 B 到 \(\mathbf{M}_{n}(\mathrm{E})\) 的 K-代数同构。令 \(d = [D : K]\) 且 \(e = [E : K]\)。我们于是有关系式

\[
[ \mathrm{A}: \mathrm{K} ] = m ^ {2} d, \qquad [ \mathrm{B}: \mathrm{K} ] = n ^ {2} e, \qquad [ \mathrm{B}: \mathrm{A} ] _ {s} = \frac {n ^ {2} e}{m ^ {2} d}
\]

以及，由公式 (30) 与 (32)，关系式

\[
i (f) = \frac {n}{m}, \qquad h (f) = \frac {n e}{m d}.
\]

当域 \(\mathrm{K}\) 是代数闭的时，我们有 \(d = e = 1\)，因此 \(i(f) = h(f)\) 且 \([\mathrm{B}:\mathrm{A}]_s = i(f)^2\)。

设 A、B 与 C 是单环，且设 \(f: A \to B\) 与 \(g: B \to C\) 是同态。设 S 是一个单 A-模。C-模 \((g \circ f)^{*}(S)\) 与 \(g^{*}(f^{*}(S))\) 同构；因此，由公式 (26) 与 (28)，我们有

\[
i (g \circ f) = \operatorname{long} _ {\mathrm{C}} \left(g ^ {*} \left(f ^ {*} (\mathrm{S})\right)\right) = i (g) \operatorname{long} _ {\mathrm{B}} \left(f ^ {*} (\mathrm{S})\right) = i (g) i (f).
\]

我们可以类似地证明等式 \(h(g \circ f) = h(g)h(f)\)。当 A 是 B 的子环且 B 是 C 的子环，并且对 f 与 g 取典范单射时，这些等式给出

\[
i (\mathrm{C}, \mathrm{A}) = i (\mathrm{C}, \mathrm{B}) i (\mathrm{B}, \mathrm{A}), \qquad h (\mathrm{C}, \mathrm{A}) = h (\mathrm{C}, \mathrm{B}) h (\mathrm{B}, \mathrm{A}).\tag{34}
\]

命题 12. —— 设 B 是一个单环，A 是 B 的一个单子环。假设 B 是一个有限生成左 A-模。设 M 是一个非零的有限生成左 B-模。令 \(A' = \text{End}_{A}(M)\) 且 \(B' = \text{End}_{B}(M)\)。则 \(B'\) 是 \(A'\) 的一个子环，环 \(A'\) 与 \(B'\) 都是单的，\(A'\) 是一个有限生成左 \(B'\)-模，并且我们有等式

\[
i (\mathrm{A} ^ {\prime}, \mathrm{B} ^ {\prime}) = h (\mathrm{B}, \mathrm{A}), \quad h (\mathrm{A} ^ {\prime}, \mathrm{B} ^ {\prime}) = i (\mathrm{B}, \mathrm{A}), \quad [ \mathrm{A} ^ {\prime}: \mathrm{B} ^ {\prime} ] _ {s} = [ \mathrm{B}: \mathrm{A} ] _ {s}.
\]

由 VIII, 第 123 页的命题 4，环 \(\mathrm{A}'\) 是单的，\(\mathrm{M}\) 是一个有限长度的 \(\mathrm{A}'\)-模，并且我们有

\[
\mathrm{long} _ {\mathrm{A}} (\mathrm{M}) = \mathrm{long} (\mathrm{A} ^ {\prime}), \quad \mathrm{long} _ {\mathrm{A} ^ {\prime}} (\mathrm{M}) = \mathrm{long} (\mathrm{A}).
\]

出于同样的理由，环 \(B'\) 是单的，并且我们有

\[
\operatorname{long} _ {\mathrm{B}} (\mathrm{M}) = \operatorname{long} \left(\mathrm{B} ^ {\prime}\right), \quad \operatorname{long} _ {\mathrm{B} ^ {\prime}} (\mathrm{M}) = \operatorname{long} (\mathrm{B}).
\]

于是，由公式 (31) 与 (30)，我们有

\[
h (\mathrm{B}, \mathrm{A}) = \operatorname{long} _ {\mathrm{A}} (\mathrm{M}) / \operatorname{long} _ {\mathrm{B}} (\mathrm{M}) = \operatorname{long} \left(\mathrm{A} ^ {\prime}\right) / \operatorname{long} \left(\mathrm{B} ^ {\prime}\right) = i \left(\mathrm{A} ^ {\prime}, \mathrm{B} ^ {\prime}\right);
\]

等式 \(h(\mathrm{A}^{\prime},\mathrm{B}^{\prime}) = i(\mathrm{B},\mathrm{A})\) 可以类似地建立。由此我们导出

\[
[ \mathrm{A} ^ {\prime}: \mathrm{B} ^ {\prime} ] _ {s} = i (\mathrm{A} ^ {\prime}, \mathrm{B} ^ {\prime}) h (\mathrm{A} ^ {\prime}, \mathrm{B} ^ {\prime}) = h (\mathrm{B}, \mathrm{A}) i (\mathrm{B}, \mathrm{A}) = [ \mathrm{B}: \mathrm{A} ] _ {s}
\]

（利用公式 (32)）。特别地，\(A'\) 是一个有限生成左 \(B'\)-模。

